% =============================================================================
%  Notas de clase — Sesión 18: Taller 2 — Resolución de problemas
%  Programación Avanzada — Universidad Panamericana
%  Compilar: pdflatex sesion18_taller_2.tex
% =============================================================================
\documentclass[12pt, oneside]{article}
\usepackage{progavanzada}
\usepackage{array}

\begin{document}

\portada{Sesión 18}{Taller 2 — Resolución de problemas}{2026}

\tableofcontents
\newpage

% =============================================================================
\section{Instrucciones del taller}
% =============================================================================

Esta sesión sigue la misma dinámica que el Taller 1: trabajo en
\textbf{equipos de dos o tres personas}, con dos momentos por problema.

\begin{description}
  \item[Parte 1 — Revisión guiada.]
    El profesor trabaja el caso de prueba en el pizarrón.
    Tu tarea es \textbf{poner atención}: observa cómo se leen los datos,
    qué operaciones se hacen y cómo se obtiene la salida esperada.
    No abras el cuaderno todavía.

  \item[Parte 2 — Trabajo en equipo.]
    Con lo observado, tu equipo generaliza el algoritmo, lo implementa en
    C++ y lo envía al juez.
\end{description}

\begin{advertencia}
  Durante la \textbf{Parte 1} no escribas código ni intentes adelantarte
  a la solución.  El objetivo es entender \emph{qué} pide el problema,
  no todavía \emph{cómo} resolverlo para cualquier entrada.
\end{advertencia}

\begin{consejo}
  \textbf{Flujo de la Parte 2:}
  \begin{enumerate}
    \item Con el ejemplo ya resuelto, identifica la condición o proceso
          general.
    \item Escribe el algoritmo en pseudocódigo antes de escribir código.
    \item Implementa en C++ y prueba con el mismo ejemplo del enunciado.
    \item Envía al juez.  Si recibes \texttt{WA}, construye un
          contraejemplo antes de modificar el código al azar.
  \end{enumerate}
\end{consejo}

\newpage

% =============================================================================
\section{Problema 1: El profesor y los nombres}
% =============================================================================

\subsection{Enunciado}

Un profesor se sabe de memoria solo algunos nombres de sus alumnos.
Cuando necesita llamar a un alumno cuyo nombre no conoce, lo hace
refiriéndose al alumno conocido más cercano.

Dados los $n$ alumnos en fila (algunos con nombre conocido, otros como
\texttt{?}), y $m$ preguntas de la forma ``¿cómo llamarías al alumno en
la posición $p$?'', escribe el programa que responda cada pregunta.

\textbf{Reglas para nombrar:}
\begin{itemize}
  \item Si la posición $p$ tiene nombre conocido: imprimir el nombre.
  \item Si el alumno conocido más cercano está a distancia $d$ y está
        a la \textbf{izquierda}: imprimir \textit{d} veces
        ``\texttt{derecha}'' seguido de ``\texttt{de Nombre}''.
  \item Si está a la \textbf{derecha}: imprimir \textit{d} veces
        ``\texttt{izquierda}'' seguido de ``\texttt{de Nombre}''.
  \item Si hay dos alumnos conocidos a la \textbf{misma distancia}
        (uno a cada lado): imprimir
        ``\texttt{justo enmedio de NombreIzq y NombreDer}''.
\end{itemize}

\textbf{Entrada.}
Primera línea: $n$ ($1 \le n \le 100$).
Segunda línea: $n$ nombres separados por espacios (\texttt{?} si se
desconoce; cadena de máximo 3 letras si se conoce).
Al menos un nombre no es \texttt{?}.
Tercera línea: $m$ ($1 \le m \le 100$).
Las siguientes $m$ líneas: una posición $p$ ($1 \le p \le n$) por línea.

\textbf{Salida.}
Por cada consulta, una línea con la forma de llamar al alumno.

\subsection{Parte 1 — Revisión guiada}

\begin{advertencia}
  \textbf{Pon atención al pizarrón.}  El profesor va a trazar la fila,
  marcar los nombres conocidos y resolver cada consulta paso a paso.
  Observa cómo se determina el nombre conocido más cercano y cómo se
  construye la cadena de respuesta.
\end{advertencia}

\begin{center}
\begin{tabular}{ll}
  \toprule
  \textbf{Entrada} & \textbf{Salida} \\
  \midrule
  \begin{minipage}[t]{5cm}
    \texttt{10}\\
    \texttt{A ? ? B ? ? ? C ? ?}\\
    \texttt{4}\\
    \texttt{3}\\
    \texttt{8}\\
    \texttt{6}\\
    \texttt{10}
  \end{minipage}
  &
  \begin{minipage}[t]{5cm}
    \texttt{izquierda de B}\\
    \texttt{C}\\
    \texttt{justo enmedio de B y C}\\
    \texttt{derecha derecha de C}
  \end{minipage} \\
  \bottomrule
\end{tabular}
\end{center}

\textbf{Notas de lo observado en el pizarrón:}

\vspace{3.5cm}

\subsection{Parte 2 — Trabajo en equipo}

\subsubsection{Análisis}

¿Cómo encuentras el nombre conocido más cercano \textbf{a la izquierda}
de la posición $p$?

\vspace{2.5cm}

¿Cómo encuentras el más cercano \textbf{a la derecha}?

\vspace{2.5cm}

¿Cuándo se imprime ``\texttt{justo enmedio}''?  ¿Cuándo sólo
``\texttt{izquierda}'' o ``\texttt{derecha}''?

\vspace{2cm}

¿Cómo generas la cadena con $d$ repeticiones de
``\texttt{derecha}'' o ``\texttt{izquierda}'' en C++?

\vspace{2cm}

\begin{recuerda}
  Las posiciones en el enunciado van de 1 a $n$.  En un arreglo de C++
  los índices van de 0 a $n-1$.  Recuerda restar 1 al leer $p$.
\end{recuerda}

\subsubsection{Pseudocódigo}

\vspace{5cm}

\newpage

% =============================================================================
\section{Problema 2: Once Cifras}
% =============================================================================

\subsection{Enunciado}

Se tiene un número de \textbf{11 cifras}.  Se parte en \textbf{tres
partes} (cada parte es una subcadena contigua de dígitos) y se suman.
¿Cuál es la menor suma posible de entre todas las formas de hacer
el corte?

\textbf{Entrada.}
Un número entero de 11 dígitos.

\textbf{Salida.}
La menor suma posible.

\subsection{Parte 1 — Revisión guiada}

\begin{advertencia}
  \textbf{Pon atención al pizarrón.}  El profesor va a mostrar qué
  significa ``partir en tres partes'', cómo se calcula la suma de
  una partición concreta y cuántas particiones distintas hay en total.
  Observa el patrón antes de pensar en la solución general.
\end{advertencia}

\begin{center}
\begin{tabular}{ll}
  \toprule
  \textbf{Entrada} & \textbf{Salida} \\
  \midrule
  \texttt{31124510988} & \texttt{3750} \\
  \bottomrule
\end{tabular}
\end{center}

\textbf{Notas de lo observado en el pizarrón:}

\vspace{3.5cm}

\subsection{Parte 2 — Trabajo en equipo}

\subsubsection{Análisis}

¿Cuántos puntos de corte necesitas para dividir una cadena de 11
dígitos en 3 partes?

\vspace{1.5cm}

Si la cadena tiene 11 dígitos y los dos puntos de corte están en las
posiciones $i$ y $j$ (con $1 \le i < j \le 10$), ¿cuántas combinaciones
$(i, j)$ hay?

\vspace{1.5cm}

¿Qué ocurre con los ceros a la izquierda al convertir una subcadena a
número?  (\textit{Ejemplo: ``\texttt{0988}'' como entero es 988.})
¿Eso afecta el resultado?

\vspace{2cm}

¿Cabe el valor máximo posible de la suma en un \texttt{int}?
(\textit{Pista: el peor caso es partir ``\texttt{99999999999}''
en ``\texttt{9}'' + ``\texttt{9}'' + ``\texttt{999999999}''.})

\vspace{2cm}

\begin{recuerda}
  Si lees el número como \texttt{string}, puedes usar \texttt{substr}
  para obtener cada parte y \texttt{stoll} para convertirla a entero.
  Iterar sobre todos los pares $(i, j)$ con dos \texttt{for} anidados
  es suficiente dado el tamaño del problema.
\end{recuerda}

\subsubsection{Pseudocódigo}

\vspace{5cm}

\newpage

% =============================================================================
\section{Resumen de la sesión}
% =============================================================================

\begin{itemize}
  \item Para ``El profesor y los nombres'': busca el nombre conocido más
        cercano a la izquierda y a la derecha; la distancia determina
        cuántas veces se repite ``\texttt{izquierda}'' o
        ``\texttt{derecha}''.
  \item Para ``Once Cifras'': hay $\binom{10}{2} = 45$ formas de partir
        11 dígitos en 3 partes; iterar sobre todas y quedarse con la
        suma mínima es la estrategia correcta.
  \item Los ceros iniciales en una subcadena no son un caso especial:
        \texttt{stoll("0988")} devuelve 988 automáticamente.
\end{itemize}

\end{document}
